\documentclass[10pt,a4paper]{amsart}
\usepackage[margin=19mm]{geometry}
\usepackage{amsmath,amssymb,mathtools}
\usepackage[scr=rsfs,cal=euler]{mathalfa}
\usepackage{xfrac}
\usepackage[unicode,hidelinks,bookmarks=false]{hyperref}
\newcommand{\ie}{i.e.\ }
\newcommand{\cf}{cf.\ }
\newcommand{\resp}{respectively\ }
\newcommand{\Subsection}{\subsection}
\providecommand{\nobreakdash}{\nobreak\hbox{-}}
\setlength{\emergencystretch}{2em}
\multlinegap0pt
\usepackage{bm}
\usepackage{bbm}
\usepackage{dsfont}
\usepackage{xspace}
\usepackage{cancel}
\usepackage{mathtools}
\usepackage{amsthm}
\makeatletter
\@ifundefined{thm}{\newtheorem{thm}{Thm}}{}
\@ifundefined{prop}{\newtheorem{prop}[thm]{Proposition}}{}
\makeatother
\newcommand{\Borbitspace}{B_{n-1}\backslash\B_{n}}
\newcommand{\B}{\mathcal{B}}
\newcommand{\F}{\mathcal{F}}
\newcommand{\G}{\mathcal{G}}
\newcommand{\Sh}{\mathcal{S}h}
\newcommand{\Sp}{\mathcal{S}p}
\newcommand{\he}{\hat{e}}
\title[$B\_{n-1}$-orbits on the flag variety and the Bruhat graph of]{$B\_{n-1}$-orbits on the flag variety and the Bruhat graph of the symmetric group}
\author{Mark Colarusso, Sam Evens}
\date{Source: arXiv:2406.10122v2 (CC BY 4.0)}
\begin{document}
\maketitle

\noindent\emph{Source and licence.} $B_{n-1}$-orbits on the flag variety and the Bruhat graph of the symmetric group. Version arXiv:2406.10122v2 (CC BY 4.0), distributed under the Creative Commons Attribution 4.0 International licence, as recorded in the arXiv licence metadata for this exact version and shown on the abstract page https://arxiv.org/abs/2406.10122v2. Full rights notice: licenses/arxiv-2406.10122-CC-BY-4.0.txt.\par\smallskip

\noindent\textit{Editorial scope.} Counted unit: the Prop whose source label is \texttt{p:Shpairs} in the article (source0.tex lines 887--894, source label \texttt{p:Shpairs}), delivered together with the article's own complete original proof of it (source0.tex lines 895--938, one \texttt{proof} environment of 44 lines). No mathematics is written, repaired or re-derived: statement and proof are the article's own text line for line. Required context retained uncounted: p:orbits (lines 759--784); eq:hatvectorF (lines 768--770); eq:tildeF (lines 773--775); eq:starF (lines 778--780). Every definition, display and label of those lines is retained unchanged. Omitted from this package and not redistributed: 1--758, 771--772, 776--777, 781--886, 939--1742. Nothing inside a retained excerpt is omitted or altered: every retained body is delivered exactly as the article's lines are written, so no omission receipt of any kind accompanies this package. Citations: the retained excerpts carry no citation command at all, so no bibliography entry of the article is transcribed and no reference is redistributed. Reference adaptation: none was needed and none was made; every reference inside the delivered unit resolves inside it, so no unresolved reference or citation is produced. Printed slips kept literal: no printed word, symbol, hypothesis or number of the counted statement or of its proof was altered. Layout and class: the article's own class is \texttt{amsart} with packages including amsmath, amssymb, mathdots, tikz, adjustbox, xy; the wrapper is the lane's pinned \texttt{amsart} editorial wrapper at 10pt on A4, which restates the theorem-like environments the retained text needs and adds the article's own macro definitions that the retained text needs (\texttt{\textbackslash B}, \texttt{\textbackslash Borbitspace}, \texttt{\textbackslash F}, \texttt{\textbackslash G}, \texttt{\textbackslash Sh}, \texttt{\textbackslash Sp}, \texttt{\textbackslash he}), with extra wrapper packages (bm, bbm, dsfont, xspace, cancel, mathtools). The delivered numbering is the unit's own. Assets: no retained span contains a figure environment, a graphics-inclusion command, a PDF-page inclusion, a table or a TikZ picture; no image file is copied into this package, the package declares no asset dependency and no image file is redistributed with it. Licence: version arXiv:2406.10122v2 of this article is distributed under the Creative Commons Attribution 4.0 International licence, as recorded in the arXiv licence metadata for this exact version and shown on the abstract page https://arxiv.org/abs/2406.10122v2. A full rights notice is delivered at licenses/arxiv-2406.10122-CC-BY-4.0.txt. Characters: every retained excerpt is pure ASCII, so the delivered engine, XeLaTeX with its default Latin Modern text font, reports no missing character.
\begin{prop}\label{p:orbits}
Let $\F\subset \B_{n}$ be a flag in standard form with 
$$
\F=(v_{1}\subset v_{2}\subset \dots\subset v_{p}\subset \dots\subset v_{n}).
$$  
(1) If $\F$ contains no hat vectors, then 
$\F=\tilde{\F}=\F^{*}$. 

\noindent (2) If $\F$ has hat vectors, we may assume that $\F$ has the form:
\begin{equation}\label{eq:hatvectorF}
\F=(v_{1}\subset\dots\subset v_{i_{k}-1}\subset\underbrace{ \he_{j_{k}}}_{i_{k}}\subset v_{i_{k}+1}\subset \dots \subset \underbrace{\he_{j_{k-1}}}_{i_{k-1}}\subset \dots \subset \underbrace{\he_{j_{1}}}_{i_{1}}\subset \dots\subset \underbrace{e_{n}}_{p}\subset v_{p+1}\subset \dots\subset v_{n}),
\end{equation}
with $j_{k}>j_{k-1}>\dots>j_{1}$ and $v_{m}$ a standard basis vector.  
Then
\begin{equation}\label{eq:tildeF}
\tilde{\F}=(v_{1}\subset\dots\subset v_{i_{k}-1}\subset\underbrace{e_{n}}_{i_{k}}\subset v_{i_{k}+1}\subset \dots\subset \underbrace{e_{j_{k}}}_{i_{k-1}}\subset\dots\subset \underbrace{e_{j_{2}}}_{i_{1}}\subset \dots\subset\underbrace{e_{j_{1}}}_{p}\subset v_{p+1}\subset\dots\subset v_{n}),
\end{equation}
%(THIS IS NOT THE MOST COMPACT WAY TO DEFINE $\tilde{\F}$ OF COURSE, BUT IT HELPS ME VISUALIZE IT.)
and 
\begin{equation}\label{eq:starF}
\F^{*}=(v_{1}\subset \dots\subset v_{i_{k-1}}\subset \underbrace{e_{j_{k}}}_{i_{k}}\subset v_{i_{k}+1}\subset \dots\subset \underbrace{e_{j_{k-1}}}_{i_{k-1}}\subset \dots\subset \underbrace{e_{j_{1}}}_{i_{1}}\subset \dots\subset\underbrace{e_{n}}_{p}\subset v_{p+1}\subset \dots\subset v_{n}),
\end{equation}
where the $v_{m}$ are the same vectors that appear in the flag in Equation (\ref{eq:hatvectorF}).
%(i.e. $\F^{*}$ is obtained from $\F$ by removing the hats). 

\end{prop}

\begin{prop}\label{p:Shpairs}
 A pair of Weyl group elements 
$(w, y)\in W\times W$ is a Shareshian pair if and only if there exists a subsequence $\Delta = \{j_{1}<j_{2}<\dots<j_{k}<n\}$ of $\{1,\dots, n\}$ containing $n$ such that: 
\begin{enumerate}
\item $y=\tau_{\Delta}w\sigma^{-1}$ and
\item $\Delta$ is a decreasing sequence for $w^{-1}$, i.e., $w^{-1}(n)<w^{-1}(j_{k})<\dots <w^{-1}(j_{1}).$
\end{enumerate}
\end{prop}

\begin{proof}
Suppose first that $Q\in\Borbitspace$.  We claim that $\Sh(Q)=(w,u^{*})$ satisfies conditions (1) and (2) of the
Proposition.  Let $\F\in Q$ be the unique flag in standard form in the orbit $Q$.  First, suppose that 
$\F$ contains no hat vectors, so that $\F=\tilde{\F}=\F^{*}$ by (1) of Proposition \ref{p:orbits}.  Let $\Delta=\{n\}$, so $\tau_{\Delta}=id$.  Since $\sigma^{-1}(\mathcal{E}^{*})=\mathcal{E_{+}}$ and $\tilde{\F}=w(\mathcal{E}_{+})$ and $\F^{*}=u^{*}(\mathcal{E}^{*})$, it follows that $u^{*}=w\sigma^{-1}$, and condition (2) is automatic.

Now suppose that $\F$ contains hat vectors.  Since $\F$ is in standard form, it is  given by (\ref{eq:hatvectorF}), so that $\tilde{\F}$ and $\F^{*}$ are given by Equations (\ref{eq:tildeF}) and (\ref{eq:starF}) 
respectively.  Let $\Delta:=\{j_{1}<j_{2}<\dots<j_{k}<n\}$, so that $\Delta$ consists of the indices of the hat 
vectors in $\F$ along with $n$.  
%are the indices of the hat vectors in $\F$.  Comparing
From Equations (\ref{eq:tildeF}) and (\ref{eq:starF}), we see that $$\F^{*}=u^{*}(\mathcal{E}^{*})=\tau_{\Delta}\tilde{\F}=\tau_{\Delta} w (\mathcal{E}_{+}).$$  
Since $\sigma^{-1}(\mathcal{E}^{*})=\mathcal{E}_{+}$, it follows that $u^{*}=\tau_{\Delta}w\sigma^{-1}$.   From (\ref{eq:tildeF}), we see $w(i_{k})=n,\, w(i_{s})=j_{s+1}$ for $s=1,\dots, k-1$, and $w(p)=j_{1}$ with $i_k < i_{k-1} < \dots < i_1 < p$, so condition (2) follows.

%by Equation (\ref{eq:tildeF}). 

Conversely, suppose that $(w,y)\in W\times W$ satisfies the conditions of (1) and (2) in the statement of the Proposition.  
Let $\Delta=\{j_{1}<\dots< j_{k}<n\}$.  By condition (2), we have $w^{-1}(n)<w^{-1}(j_{k})<\dots< w^{-1}(j_{1})$.  
Let $w^{-1}(n)=i_{k}$, $w^{-1}(j_{k})=i_{k-1}, \dots w^{-1}(j_{2})=i_{1}$, and $w^{-1}(j_{1})=p$.  Let $\mathcal{G}:=w(\mathcal{E}_{+})$.  
Then 
\begin{equation*}
\begin{split}
&\mathcal{G}=
(e_{w(1)}\subset\dots\subset e_{w(i_{k}-1)}\subset\underbrace{e_{n}}_{i_{k}}\subset e_{w(i_{k}+1)}\subset \dots\subset \underbrace{e_{j_{k}}}_{i_{k-1}}\subset\dots\subset \underbrace{e_{j_{2}}}_{i_{1}}\subset \dots\subset\underbrace{e_{j_{1}}}_{p}\subset\\
&\subset e_{w(p+1)}\subset\dots\subset e_{w(n)}),
\end{split}
\end{equation*}
and it follows that 
\begin{equation*}
\begin{split}
&\tau_{\Delta} (\mathcal{G})=(e_{w(1)}\subset \dots\subset e_{w(i_{k-1})}\subset \underbrace{e_{j_{k}}}_{i_{k}}\subset e_{w(i_{k}+1)}\subset \dots\subset \underbrace{e_{j_{k-1}}}_{i_{k-1}}\subset \dots\subset \underbrace{e_{j_{1}}}_{i_{1}}\subset \dots\subset\underbrace{e_{n}}_{p}\subset\\ &\subset e_{w(p+1)} 
\subset \dots\subset e_{w(n)}).
\end{split}
\end{equation*}
Now let $\F$ be the flag in standard form:
\begin{equation*}
\begin{split}
&\F=(e_{w(1)}\subset\dots\subset e_{w(i_{k}-1)}\subset\underbrace{ \he_{j_{k}}}_{i_{k}}\subset e_{w(i_{k}+1)}\subset \dots \subset \underbrace{\he_{j_{k-1}}}_{i_{k-1}}\subset \dots \subset \underbrace{\he_{j_{1}}}_{i_{1}}\subset \dots\subset \underbrace{e_{n}}_{p}\subset\\
&\subset e_{w(p+1)}\subset \dots\subset e_{w(n)}),
\end{split}
\end{equation*}
Then it follows from Proposition \ref{p:orbits} that $\tilde{\F}=\mathcal{G}$ and $\F^{*}=\tau_{\Delta}(\G)$.  
But $\G=w(\mathcal{E}_{+})$ and $$\tau_{\Delta}(\G)=\tau_{\Delta}w(\mathcal{E}_{+})=\tau_{\Delta}w\sigma^{-1}(\mathcal{E}^{*}).$$  
Thus, $\Sh(B_{n-1}\cdot \F)=(w, \tau_{\Delta}w\sigma^{-1})=(w,y)$ by condition (1).  
Thus, $(w,y)\in \Sp$ as asserted. 
\end{proof}

\end{document}
